\documentclass[11pt]{article}

\usepackage[margin=1in]{geometry}
\usepackage{amsmath,amssymb,amsthm}
\usepackage{graphicx}
\usepackage[hidelinks]{hyperref}

\newtheorem{theorem}{Theorem}
\newtheorem{lemma}{Lemma}

\newcommand{\T}{\mathbb T}
\newcommand{\D}{\mathbb D}
\newcommand{\N}{\mathbb N}
\newcommand{\cC}{\mathcal C}
\newcommand{\norm}[1]{\left\lVert #1 \right\rVert}

\title{A Counterexample to the Clark Cauchy Transform Converse}
\author{Automated Math Discovery Smoke Test}
\date{July 3, 2026}

\begin{document}
\maketitle

\section*{Source and Open Problem}

This packet concerns Open Problem 4.4 from:

\begin{quote}\raggedright
C. Bellavita, E. Dellepiane, A. Hartmann, and J. Mashreghi,
\emph{Infinitely supported harmonically weighted Dirichlet spaces which are
de Branges Rovnyak spaces}, arXiv:2509.04907.
\end{quote}

The source proves that if $u$ is a one-component inner function and
$\sigma$ is its Clark measure, then the truncated Cauchy transform
\[
  \cC_\sigma f(\zeta)
  =
  \int_{\T\setminus\{\zeta\}}
    \frac{f(s)}{1-\overline{s}\zeta}\,d\sigma(s)
\]
is bounded on $L^2(\sigma)$. It then asks whether the converse is true:

\begin{quote}
Open problem 4.4. Let $\sigma$ be the Clark measure associated with $u$.
Is it true that if $\cC_\sigma$ is bounded on $L^2(\sigma)$, then $u$ is a
one-component inner function?
\end{quote}

\begin{center}
\includegraphics[width=\textwidth]{figures/open_problem_crop.png}
\end{center}

\section*{Idea of the Proof}

Boundedness of the full Cauchy transform controls the singular integral
operator, but it does not force the geometric mass-gap conditions appearing
in Bessonov's characterization of Clark measures of one-component inner
functions. We make the atom masses so small that the Cauchy matrix is
Hilbert-Schmidt. At the same time, the masses are far too small compared to
the adjacent gaps, so Bessonov's necessary lower estimate fails.

\section*{The Counterexample}

\begin{theorem}
There is an inner function $u$ with Clark measure $\sigma$ such that
$\cC_\sigma$ is bounded on $L^2(\sigma)$, but $u$ is not one-component.
\end{theorem}

\begin{proof}
For $n\ge1$ put
\[
  \theta_n=2^{-n},\qquad
  \zeta_n=e^{i\theta_n},\qquad
  m_n=2^{-4n},
\]
and define the finite positive singular measure
\[
  \sigma=\sum_{n=1}^{\infty} m_n\delta_{\zeta_n}.
\]
By the standard Herglotz--Clark correspondence for singular measures on
$\T$, there is an inner function $u$ for which $\sigma$ is the Clark measure
corresponding to the parameter $1$.

We first prove that $\cC_\sigma$ is bounded on $L^2(\sigma)$. Identifying
$L^2(\sigma)$ with $\ell^2(m_n)$, the transform is represented on finitely
supported sequences by
\[
  (\cC_\sigma f)(\zeta_j)
  =
  \sum_{i\ne j}
    \frac{m_i f(\zeta_i)}
         {1-\overline{\zeta_i}\zeta_j}.
\]
Conjugating by the unitary map
$f\mapsto (m_n^{1/2}f(\zeta_n))_{n\ge1}$ gives the matrix
\[
  a_{ji}
  =
  \begin{cases}
  \displaystyle
  \frac{(m_i m_j)^{1/2}}
       {1-\overline{\zeta_i}\zeta_j},& i\ne j,\\[1.2ex]
  0,&i=j.
  \end{cases}
\]
It is enough to show that this matrix is Hilbert-Schmidt.

For $i<j$,
\[
  |\theta_i-\theta_j|
  =
  2^{-i}-2^{-j}
  \ge 2^{-i-1}.
\]
Since $|\theta_i-\theta_j|<1$, the elementary estimate
$|1-e^{it}|\ge c|t|$ gives
\[
  |1-\overline{\zeta_i}\zeta_j|
  =
  |1-e^{i(\theta_j-\theta_i)}|
  \ge c\,2^{-i}.
\]
Therefore
\[
  |a_{ji}|^2
  \le
  C\,2^{-4i}2^{-4j}2^{2i}
  =
  C\,2^{-2i-4j}
  \qquad (i<j).
\]
The same estimate with $i$ and $j$ interchanged handles $j<i$. Hence
\[
  \sum_{i,j\ge1}|a_{ji}|^2
  \le
  2C\sum_{1\le i<j}2^{-2i-4j}
  <\infty.
\]
Thus $\cC_\sigma$ is Hilbert-Schmidt, in particular bounded, on
$L^2(\sigma)$.

It remains to show that $u$ is not one-component. The source quotes
Bessonov's characterization as Theorem 3.8. One necessary condition for a
Clark measure of a one-component inner function is the existence of a
constant $A>0$ such that every isolated atom $\xi$ satisfies
\[
  \sigma(\{\xi\})\ge A|\xi-\xi^+|
\]
for each adjacent atom $\xi^+$.
Here $\zeta_{n+1}$ is adjacent to $\zeta_n$ in the accumulating chain, and
\[
  |\zeta_n-\zeta_{n+1}|
  \asymp
  |\theta_n-\theta_{n+1}|
  =
  2^{-n-1}.
\]
Consequently
\[
  \frac{\sigma(\{\zeta_n\})}{|\zeta_n-\zeta_{n+1}|}
  \asymp
  \frac{2^{-4n}}{2^{-n}}
  =
  2^{-3n}
  \longrightarrow 0.
\]
No positive lower mass-gap constant $A$ exists. Therefore Bessonov's
necessary condition fails, and the inner function $u$ cannot be
one-component.
\end{proof}

\section*{Verification and Novelty Check}

The file \texttt{code/check\_hilbert\_schmidt.py} computes finite partial
sums of the Hilbert-Schmidt norm for the displayed matrix. This is only a
sanity check; the proof above gives the exact summability estimate.

Before promotion, the run indexes were searched for \texttt{2509.04907},
the source title, \texttt{Clark measure}, \texttt{Cauchy transform}, and
\texttt{one-component inner function}. Web searches for exact and nearby
phrases from Open Problem 4.4 only found the source paper. The counterexample
uses the problem exactly as stated; it does not address possible strengthened
versions where one assumes Bessonov's atom-neighbor geometry or mass-gap
conditions in advance.

\section*{References}

\begin{thebibliography}{9}

\bibitem{BDHM}
C. Bellavita, E. Dellepiane, A. Hartmann, and J. Mashreghi,
\emph{Infinitely supported harmonically weighted Dirichlet spaces which are
de Branges Rovnyak spaces}, arXiv:2509.04907, 2025.

\bibitem{Bessonov}
R. V. Bessonov,
\emph{Duality theorems for coinvariant subspaces of $H^1$},
Advances in Mathematics \textbf{271} (2015), 62--90.

\bibitem{CMR}
J. A. Cima, A. L. Matheson, and W. T. Ross,
\emph{The Cauchy Transform},
Mathematical Surveys and Monographs, American Mathematical Society, 2006.

\end{thebibliography}

\end{document}
